\textbf{Infeasible region.} The main limitation is physical rather than algorithmic. By \eqref{eq:feas}, no quasi-static release from rest can take the rod to a shallow target without passing it once $r$ exceeds $\varphi_{\lim}/\sin\varphi_{\lim}$. For $r=1.1$ this rules out every target shallower than about 0.75\,rad. In this region the shield can only defer, which reproduces the model controller's overshoot, or decline, which stops short. A fast slip-triggered re-clamp (a ``catch'') could reach these targets: in a separate probe, a 50\,ms detection-plus-clamp delay limited the first-slip jump to 0.03--0.07\,rad even at $r=1.5$. However, it needs slip detection far faster than our 10\,s current-based angle measurement, and we leave it to future work.

\textbf{Certified under a model, not under the world.} Proposition~\ref{prop:pivot} assumes a rigid rod, torsional capacity linear in grip, constant kinetic friction and no stored pad energy. In simulation the bound held on average, with simulated travel at 0.90--0.97 of the bound, but 24 of 128 direct test releases exceeded it by up to 0.048\,rad. We attribute this to soft-contact geometry, and it exceeds our 0.02\,rad sensing margin. On hardware, three effects could make the bound optimistic: soft-finger friction grows sub-linearly with load \cite{xydas1999}, so capacity falls faster than grip when grip is reduced; elastomer pads store torsional energy before breakaway; and static friction rises with dwell time \cite{dieterich1972}. A practical value of $r_{\max}$ should therefore be measured on the actual pad--object pair, at the hold times the controller uses, and taken from the upper tail of the measured ratios.

\textbf{Simulator.} MuJoCo has no true stick state \cite{mujocoissue}. We emulated static and kinetic friction by switching coefficients with hysteresis. The emulation reproduced the intended ratio (an emulated 1.2 measured 1.206) and an intrinsic ratio of 1.008, and first-slip travel was insensitive to the time step and switching thresholds. It is nonetheless a step-function Stribeck curve without rate-and-state effects \cite{canudas1995}. Kolbert et al.\ found that MuJoCo under-predicted a real pivot \cite{kolbert2016}, so absolute travel values should not be transferred to hardware.

\textbf{Learning and scope.} We used linear policies, evolution strategies, 1560 episodes per run and 5 seeds per method. Larger policies or other optimisers could change the comparison between shield-in-training and shield-at-test. Sensing during training was a noise model, not the motor-current pipeline. Lift-with-slip, the third task we originally planned, is not studied.
